% Phase 2 methodology (Harshil)
% Code: src/dl_data.py (loader, normalisation), src/dl_models.py (CNN, U-Net), src/dl_losses.py,
% src/train_dl.py (training loop), src/dl_gbm_oof.py (hybrid input), src/evaluate.py (metrics, shared with
% Phase 1). Parameter counts from build_model; numbers from reports/dl/results_dl.csv.

\section{Deep Learning Methodology}
\label{sec:dl_method}

Phase~2 keeps the task, the data, the splits and the metrics of Phase~1 (Section~\ref{sec:method}) and
replaces the per-pixel model with a network that maps a whole patch to a whole probability map. Everything
that could make the two phases differ for reasons other than the model is held fixed: the same official
split, the same masking of label $-1$, the same \texttt{evaluate()} function, and thresholds chosen on
eval only.

\subsection{Data loader and normalisation}

Each sample is the full $64\times64$ patch: an input tensor of shape $(12, 64, 64)$ and a label map of
shape $(64, 64)$ with values in $\{-1, 0, 1\}$. We use all 14\,979 / 1\,877 / 1\,689 train / eval / test
patches. There is no pixel subsampling, unlike the 1:5 subsampled pixel tables of Phase~1, because a
convolutional network needs every pixel of the patch as context.

The raw export contains impossible values (for example a wind direction of $-505\,870^\circ$), and the
channels differ in scale by several orders of magnitude. Each of the 11 continuous channels $c$ is
therefore clipped to its 0.5th and 99.5th percentile on the training split, $[\ell_c, u_c]$, and then
standardised:
\begin{equation}
\tilde x_c = \frac{\operatorname{clip}(x_c, \ell_c, u_c) - \mu_c}{\sigma_c},
\label{eq:dl_norm}
\end{equation}
where $\mu_c$ and $\sigma_c$ are the mean and standard deviation of the clipped training values. All
statistics come from the training split only and are applied unchanged to eval and test. The fire mask
\texttt{PrevFireMask} is already in $\{-1, 0, 1\}$ and is left as it is. Labels of $-1$ are kept in the
tensors and removed inside every loss and metric. The training loader shuffles patches with a seeded
generator, and eval and test are read in a fixed order. Augmentation (Section~\ref{sec:dl_aug}), when
used, is applied to training batches only.

\subsection{Models}

Both models map $(B, C, 64, 64)$ to $(B, 1, 64, 64)$ logits, one per pixel (Fig.~\ref{fig:dl_arch}).
\begin{itemize}
\item \textbf{CNN.} Four $3\times3$ convolutions with batch normalisation and ReLU, 32 channels, no
downsampling, and a $1\times1$ output convolution (31\,393 parameters). Its receptive field is
$9\times9$\,px, about the reach of the Phase~1 neighbourhood features, so it tests whether a learned
local filter beats engineered local features.
\item \textbf{U-Net}~\cite{ronneberger2015}. An encoder of $D$ levels: each level is a block of two
$3\times3$ convolutions with batch normalisation and ReLU, and levels are joined by $2\times2$ max
pooling. The channel count doubles at each level from a base width $W$. The decoder upsamples with
$2\times2$ transposed convolutions, concatenates the encoder features of the same resolution (skip
connections), applies a block, and ends in a $1\times1$ convolution. The ablation varies $D \in
\{2, 3, 4\}$ and $W \in \{16, 32, 64\}$. The final model has $D = 4$, $W = 32$ (1\,929\,601 parameters;
469\,761 at $D = 3$).
\end{itemize}
Both are fully convolutional, so the same weights apply at every pixel and any input size works.

\subsection{Training}

Training minimises a loss over the labelled pixels of each batch: weighted BCE, focal loss or BCE plus
Dice (Section~\ref{sec:dl_loss}). The fire weight $w = 90.0$ is the no-fire to fire ratio of the training
pixels. We use Adam~\cite{kingma2015} with learning rate $10^{-3}$, batch size 32 and at most 30 epochs.
After every epoch the model predicts the eval split, and eval PR AUC is computed over its labelled pixels.
The best epoch is saved, and training stops after 5 epochs without improvement~\cite{prechelt1998}. The
saved model then predicts eval and test. The decision threshold is the one that maximises F1 on eval, and
it is applied unchanged to test. Test scores go through the Phase~1 \texttt{evaluate()} function into a
separate results file, so Phase~1 results are never overwritten. Every run fixes the Python, NumPy and
PyTorch seeds and the loader seed. Section~\ref{sec:dl_validation} reports how much results still move
between seeds and machines (Apple-silicon GPUs and an NVIDIA T4).

\begin{table}[t]
\centering
\caption{Final Phase~2 model. Depth and loss differ from the defaults of \texttt{src/train\_dl.py}.}
\label{tab:dl_final}
\footnotesize
\renewcommand{\arraystretch}{1.1}
\begin{tabular}{l>{\raggedright\arraybackslash}p{5.9cm}}
\toprule
\rowcolor{tblhead}
\textbf{Setting} & \textbf{Value} \\
\midrule
Input & 12 channels, $64\times64$, train-only clip and z-score \\
Architecture & U-Net, 4 levels, base width 32, 1.93\,M parameters \\
Loss & focal, $\gamma = 2$, fire weight 90.0, label $-1$ masked \\
Optimiser & Adam, learning rate $10^{-3}$, batch size 32 \\
Stopping & eval PR AUC, patience 5, at most 30 epochs \\
Augmentation & none \\
Threshold & best F1 on eval \\
Seeds & 0, 1, 2 (reported as mean $\pm$ std) \\
\bottomrule
\end{tabular}
\end{table}

\subsection{Hybrid input}

The hybrid U-Net adds the Phase~1 XGBoost fire probability as a 13th input channel. The XGBoost
hyperparameters are those of Phase~1. On eval and test, the channel comes from XGBoost fitted on the whole
Phase~1 training pixel table. On the training patches, it comes from five out-of-fold models (Eq.~\ref{eq:oof}),
with folds grouped by fire event, so no fold model has seen the fire it predicts. Events are found as in
Phase~1, by linking patches that share a $16\times16$ elevation block. The channel is standardised with its
training mean and standard deviation. The probability ensembles average U-Net and XGBoost probabilities
with a weight tuned on eval.

\subsection{Experiments}

Each experiment changes one setting of the default model (U-Net, 3 levels, width 32) and keeps the
rest. The augmentation runs use focal loss, chosen by the loss sweep, and the other ablations use
weighted BCE:
\begin{itemize}
\item \textbf{baseline}: CNN against U-Net;
\item \textbf{loss}: weighted BCE, focal or BCE + Dice, 3 seeds (Section~\ref{sec:dl_loss});
\item \textbf{augmentation}: none, flips, rotations or both, with the wind direction turned, 3 seeds
(Section~\ref{sec:dl_aug});
\item \textbf{architecture}: depth 2/3/4 at width 32, and width 16/32/64 at depth 3;
\item \textbf{channel groups}: drop weather, vegetation, terrain, fire or population;
\item \textbf{crop}: train on centre $32\times32$ crops, score on full patches;
\item \textbf{hybrid}: with and without the XGBoost channel, and probability ensembles.
\end{itemize}
The architecture, channel, crop and hybrid runs use one seed, so differences below 0.02, the largest
shift between same-seed reruns, are treated as noise. The final model combines the winning settings and
is retrained on three seeds (Table~\ref{tab:dl_final}). All choices are made on eval. Test is reported
but never used to choose.
